\documentclass[letterpaper,11pt]{article}
\usepackage{graphicx}
\usepackage{geometry}
\linespread{1.2}                        
\setlength{\oddsidemargin}{0.0in}            
\setlength{\textwidth}{6.5in}    
\usepackage[utf8]{inputenc} % Asegúrate de usar utf8
\usepackage[T1]{fontenc} % Asegúrate de usar T1 font encoding      
\setlength{\topmargin}{-.6in}            
\setlength{\textheight}{9in}
\usepackage{latexsym} 
\usepackage{amssymb} 
\usepackage{amsmath}
\usepackage{amsxtra}
\usepackage[mathcal]{eucal} 
\usepackage{enumerate}
\usepackage{hyperref}
\usepackage{rotating}
\usepackage{caption}
\usepackage{lscape}
\usepackage{paralist} 
\usepackage{indentfirst}
\usepackage[dvipsnames,svgnames]{xcolor} 
\usepackage{setspace}
\usepackage{multicol} 
\usepackage{float}

\definecolor{codegreen}{rgb}{0,0.6,0}
\definecolor{codegray}{rgb}{0.5,0.5,0.5}
\definecolor{codepurple}{cmyk}{0.65, 1, 0, 0.35}
\definecolor{backcolour}{rgb}{0.95,0.95,0.92}

\usepackage{listings}
\usepackage{xcolor}

\lstset{
    language=R, % lenguaje de programación
    basicstyle=\small\ttfamily, % estilo básico
    commentstyle=\color{gray}, % estilo para los comentarios
    keywordstyle=\color{blue}, % estilo para las palabras clave
    stringstyle=\color{black}, % estilo para las cadenas de texto
    showstringspaces=false, % no mostrar espacios en cadenas de texto
    tabsize=4, % tamaño de la tabulación
    breaklines=true, % permitir el salto de línea automático
    postbreak=\mbox{\textcolor{red}{$\hookrightarrow$}\space}, % símbolo al final de una línea partida
    numbers=left, % mostrar números de línea
    numberstyle=\tiny\color{gray}, % estilo para los números de línea
    frame=single, % mostrar un borde alrededor del código
    backgroundcolor=\color{gray!10}, % color de fondo
}

\begin{document}

\begin{center}
    {\Large  Taller 2: Econometría }

Profesor:    Pedro Comte\\
Ayudantes: Jociney Honorio y Fabian Cisternas
\end{center}

\textbf{Instrucciones:} 

\begin{itemize}
    \item Usted dispone de 120 minutos  para intentar los 100 puntos del taller.
    \item Cuide la presentación y redacción de sus respuestas.
    \item Puede utilizar su computador y los apuntes tanto de clase como de ayudantía.
    \item Debe subir al aula el archivo de colab o el link (extensión .ipynb)
\end{itemize}

\section{Trabajo Base de Datos (40 puntos)}

Para responder las preguntas siguientes, utilice la base de datos \textit{ceosal1} de la librería \textbf{wooldridge}. Asegúrese de cargarla correctamente (sin nulos) y de especificar el programa de forma rigurosa para evitar cambios por aleatoriedad.

\begin{enumerate}

\item (10 puntos) Examine la base y elabore una breve descripción de sus variables, distinguiendo entre dependientes, explicativas, continuas y \textit{dummies}. \textquestiondown Para qué tipo de estudio económico se emplea habitualmente esta base de datos?



    \item (10 puntos) Genera una variable indicadora (\textit{dummy}) que tome el valor 1 cuando la variable \textit{sales} supere su promedio y 0 en caso contrario. Asigne a esta nueva variable un nombre que sea fácilmente identificable. \textquestiondown Que efecto, a priori, causaría esta variable en el salario?
    


    \item (5 puntos) Cree una variable de la multuplicación entre la variable \textit{dummy} antes creada y la variable \textit{indus}. \textquestiondown Qué indica esta interacción? 




    \item (10 puntos) Cree una variable de la multiplicación entre la variable continua \textit{roe}  y la variable \textit{dummy} antes creada en el punto 2. \textquestiondown Qué indica esta interacción? \textquestiondown Que efecto, a priori, causaría esta variable en el salario y de que forma afecta?



    \item (5 puntos) Cree una variable que indique si el compañia corresponde a industria (\textit{indus}) ó si \textit{finance}.  Asigne a esta nueva variable un nombre que sea fácilmente identificable.



\end{enumerate}

\section{Modelos de Regresión Lineal (60 puntos)}

\begin{enumerate}

    \item (20 puntos) Genere un modelo de regresión lineal, sobre alguna de las variables dependientes, con respecto a las variables \textit{sales}, \textit{roe} y \textit{ros}. Interprete sus resultados y la significancia de los regresores. \textquestiondown El modelo tiene significancia global al 5\%?



\item (20 puntos) Al modelo de regresión lineal anterior, agregue las variables creadas en el inciso 1.2 y 1.4, interprete sus resultados. \textquestiondown Cómo afectan los regresores de estas variables en el salario? ¿Son significativos al 5\% individualmente?



\item (20 puntos) Realiza un test de hipotesis para evaluar si las dos variables que se agregaron son significativas para el modelo simultaneamente e interprete su resultado. Use significancia del 5\%.


\end{enumerate}

\end{document}